\chapter{\chapterOne}

\section{Videojuegos}
\subsection{Definición y características}


Para establecer una base conceptual sólida, conviene comenzar por definir qué se entiende por juego. \textcite{stenros_game_2017} realiza una revisión de más de 60 definiciones de distintos autores desde los años 30, y propone una síntesis coherente: “Un juego es un sistema basado en reglas con un resultado variable y cuantificable, en el que se asignan diferentes valores a los distintos resultados; el jugador realiza un esfuerzo para influir en dicho resultado, se siente emocionalmente vinculado a él, y las consecuencias de la actividad son negociables”.

En cuanto al videojuego, Esposito, citado por Arjoranta \cite{arjoranta_how_2019}, propone una definición inicial: “Un videojuego es un juego que jugamos gracias a un aparato audiovisual y que puede estar basado en una historia”. Aunque minimalista, esta definición resulta útil como punto de partida. \textcite{deterding2013modes} amplía esta idea al señalar que “un videojuego existe en la medida en que las personas lo interpretan y lo viven como tal, dentro de un contexto con reglas, significados y materiales que permiten esa experiencia”. En este sentido, un videojuego puede concebirse como una prueba mental ejecutada en un entorno computacional bajo ciertas reglas, cuyo fin es la diversión, el entretenimiento o la obtención de una recompensa, emocional o no.

La característica más distintiva de los videojuegos, y la que los diferencia de otros medios narrativos tradicionales como la literatura o el cine, es la \textbf{interactividad}. Este componente constituye la esencia del medio videolúdico y le confiere su naturaleza participativa \cite{lopez_canicio_cuando_2019}. La interactividad permite que el productor diseñe una narrativa que no solo es recibida e interpretada por el jugador, sino también co-construida activamente por él a través de sus decisiones y acciones. 

\subsection{Taxonomía en los videojuegos}

Uno de los aspectos más relevantes al estudiar los videojuegos es su
\textbf{taxonomía} o clasificación, ya que permite comprender su diversidad
y las distintas formas en que los jugadores interactúan con ellos. Existen
múltiples criterios de clasificación que delimitan su alcance y
características.

La forma más simple de clasificación es según el \textit{hardware o
	plataforma} en la que se ejecutan, pudiendo ser videojuegos de consola, PC,
dispositivos móviles, portátiles, entre otros
\cite{herrero_taxonomivideojuegos_2015}. No obstante, esta tipología ha
perdido valor en la actualidad, dado que muchos títulos se desarrollan de
manera multiplataforma, disponibles simultáneamente en distintos dispositivos.

Una clasificación más precisa y ampliamente utilizada es aquella basada en
el \textit{gameplay} o estilo de juego, distinguiendo géneros como acción,
aventura, simulación, estrategia, rol (RPG) y casual, entre otros
\cite{herrero_taxonomivideojuegos_2015}. Es importante señalar que estos
géneros no existen de manera aislada, sino que con frecuencia se combinan
entre sí para dar lugar a experiencias híbridas y más complejas. Esta
tendencia demuestra que, aunque las clasificaciones son útiles para el
análisis, un videojuego raramente pertenece de forma exclusiva a un solo
género \cite{Starosta_Kiszka_Szyszka_Starzec_Strojny_2024,
	grelier2023datadrivenclassificationsvideogame}.

Para los fines de esta investigación, resulta especialmente relevante la
clasificación que distingue entre los \textbf{videojuegos tradicionales} y
los \textit{serious games}. Los primeros tienen como finalidad principal el
entretenimiento, mientras que los \textit{serious games} combinan elementos
lúdicos con fines educativos, de capacitación o sensibilización, enfatizando
el aprendizaje y el desarrollo de habilidades específicas
\cite{bontchev_serious_2015}. Dentro de esta última categoría, el subgénero
de la \textbf{novela visual} merece una atención particular: se caracteriza
por colocar la narrativa y la toma de decisiones como ejes centrales de la
experiencia, presentando un formato similar a una novela ilustrada donde el
jugador influye activamente en el desarrollo del relato
\cite{arsenault_2009, Qaffas_2020}. Esta estructura narrativa lo convierte
en un formato especialmente adecuado para contextualizar contenidos de
aprendizaje dentro de una historia significativa y cercana al estudiante,
aspecto que fundamenta su elección como género base en el diseño de la
solución propuesta.




\section{Videojuegos como herramienta educativa}

\subsection{Videojuegos como herramienta educativa}

En las últimas décadas, los videojuegos han dejado de ser considerados únicamente una forma de entretenimiento para convertirse en una herramienta con un potencial significativo en el ámbito educativo. Su capacidad para integrar elementos visuales, narrativos e interactivos les permite generar entornos de aprendizaje más atractivos, dinámicos y participativos, favoreciendo la motivación y el compromiso de los estudiantes \cite{ashinoff_potential_2014}.

Los videojuegos aplicados a la educación, se desarrollan con un propósito que trasciende la simple diversión. Estos videojuegos combinan las mecánicas propias del juego con objetivos formativos o de desarrollo de competencias, ofreciendo experiencias que permiten al usuario aprender mediante la acción, la experimentación y la toma de decisiones dentro de un entorno controlado \cite{padron_gonzalez_uso_2023}.

El valor educativo de los videojuegos radica en su capacidad para simular situaciones complejas del mundo real de forma segura y accesible. A través de estos entornos virtuales, los estudiantes pueden enfrentarse a problemas, tomar decisiones, observar las consecuencias de sus acciones y reflexionar sobre los resultados. Este tipo de aprendizaje activo y contextual favorece el desarrollo de habilidades cognitivas, estratégicas y sociales, al mismo tiempo que refuerza conocimientos específicos del área de estudio \cite{ashinoff_potential_2014,padron_gonzalez_uso_2023}.

Diversas áreas del conocimiento se han beneficiado de la aplicación de videojuegos educativos. En ciencias, por ejemplo, los juegos permiten visualizar fenómenos físicos o biológicos de difícil observación directa; en matemáticas, ayudan a desarrollar el razonamiento lógico y la resolución de problemas mediante desafíos interactivos \cite{ibrahim_students_2010,mathew_teaching_2019,leutenegger_games_nodate,silva_systematic_2020}; y en humanidades, favorecen la comprensión de contextos históricos, culturales o lingüísticos a través de la inmersión narrativa. Incluso en entornos profesionales y universitarios, los videojuegos se han utilizado para la capacitación en habilidades técnicas, toma de decisiones, gestión de recursos y liderazgo.

\section{Fundamentos psicológicos y pedagógicos del aprendizaje}

\subsection{Teoría del constructivismo}

El \textbf{constructivismo} constituye una de las corrientes más influyentes en la educación contemporánea y concibe el aprendizaje como un proceso activo de construcción de conocimiento. Desde esta perspectiva, el estudiante interpreta, reorganiza e integra nuevos contenidos a partir de sus experiencias previas y esquemas mentales, por lo que el aprendizaje trasciende la transmisión pasiva de información \cite{tigse_parreno_constructivismo_2019,ortiz_granja_constructivismo_2015}.

En coherencia con lo anterior, la enseñanza se entiende como un proceso dialógico y participativo en el que el docente asume un rol de mediador: orienta la actividad cognitiva y promueve la reflexión crítica. Además, el error se asume como un insumo formativo que permite reajustar estructuras cognitivas, favoreciendo la autonomía, la autorregulación y la transferencia del conocimiento. Asimismo, el contexto social y cultural (interacción, lenguaje y colaboración) funciona como un mediador clave del aprendizaje \cite{ortiz_granja_constructivismo_2015}.

En el ámbito de los videojuegos, estas ideas resultan especialmente pertinentes porque el jugador actúa, toma decisiones, observa consecuencias y construye significados a partir de la interacción con un entorno virtual, lo que permite interpretar el videojuego como un entorno con potencial constructivista \cite{lukic_konstruktivizam_nodate}.

\subsection{Aprendizaje basado en problemas}

El \textbf{Aprendizaje Basado en Problemas (ABP)} se inscribe en las teorías constructivistas del aprendizaje y constituye una metodología centrada en el estudiante, orientada a la resolución de situaciones complejas y contextualizadas \cite{ruiz-meza_application_nodate}.
Su objetivo fundamental es promover un aprendizaje significativo a través del análisis, la reflexión y la colaboración, permitiendo que los estudiantes construyan activamente su conocimiento al enfrentarse a problemas reales o simulados que demandan la aplicación integrada de saberes teóricos y prácticos.

En este enfoque, los problemas no se conciben como simples ejercicios de aplicación mecánica, sino como detonantes del pensamiento crítico, la formulación de hipótesis y la búsqueda de soluciones fundamentadas. De esta manera, el ABP impulsa la autonomía intelectual, la autorregulación del aprendizaje y el desarrollo de habilidades de razonamiento superior, transformando el papel del docente, que pasa de ser un transmisor de información a un facilitador o guía del proceso de descubrimiento \cite{Ge_Zhu_Lin_Jiang_Li_Lu_Mi_Tung_2025}.

Diversos estudios han demostrado la eficacia del ABP en la mejora de las competencias cognitivas y profesionales de los estudiantes. Por ejemplo, \textcite{Lu_Singh_2024} evidencia que esta metodología incrementa el pensamiento crítico y el rendimiento académico de los universitarios al situarlos en contextos de aprendizaje activo y reflexivo.
Del mismo modo, investigaciones recientes destacan la relevancia del aprendizaje colaborativo en línea, señalando que la resolución conjunta de problemas —potenciada por entornos digitales interactivos— favorece la construcción colectiva del conocimiento y el desarrollo de habilidades sociales y cognitivas de orden superior \cite{Jiang_Ruan_Feng_Jiang_Xiong_2023}.

En el contexto de la educación superior, el ABP promueve una transformación metodológica hacia modelos de enseñanza más participativos, reflexivos y orientados a la resolución de problemas reales. Esta transición contribuye a fortalecer la relación entre teoría y práctica, facilitando la adquisición de competencias transferibles y duraderas \cite{ruiz-meza_application_nodate}.


\subsection{Teoría del flujo}

El concepto de Estado de Flujo fue introducido por \textcite{Csikszentmihalyi_2008}, quien lo definió como una experiencia óptima de inmersión total, disfrute y concentración absoluta en una actividad que resulta intrínsecamente gratificante. Durante el flujo, el individuo experimenta una pérdida de la noción del tiempo, una disminución de la autoconciencia y un sentido de control total sobre sus acciones. Este fenómeno psicológico se asocia con un equilibrio preciso entre el nivel de desafío de la tarea y las habilidades del participante, lo que genera una sensación de compromiso pleno y satisfacción por la propia actividad, más allá de recompensas externas.

En el contexto educativo, el estado de flujo se ha convertido en un marco teórico clave para comprender los procesos de motivación, atención y disfrute en el aprendizaje. Según \textcite{Schmidt_2010}, cuando las actividades de aprendizaje logran mantener un equilibrio adecuado entre la dificultad y la competencia percibida del estudiante, se facilita la aparición del flujo, lo que conduce a una mayor implicación cognitiva, persistencia y rendimiento académico. En este sentido, el flujo no solo potencia la adquisición de conocimientos, sino que también promueve el desarrollo de habilidades metacognitivas, la creatividad y el bienestar emocional del estudiante \cite{Csikszentmihalyi_2008}.

La teoría del flujo ha demostrado ser especialmente relevante en el ámbito del aprendizaje mediado por juegos y entornos digitales interactivos. Los videojuegos y las experiencias gamificadas, al incorporar retos progresivos, retroalimentación inmediata y metas claras, reproducen de manera natural las condiciones que facilitan la inmersión y el compromiso sostenido. \textcite{Admiraal} evidencia que, en contextos de aprendizaje basados en juegos colaborativos, los estudiantes que alcanzan un estado de flujo tienden a mostrar mayor implicación y disfrute durante la actividad, lo que influye positivamente en su desempeño y en su conexión emocional con el proceso educativo. Aunque en algunos casos el flujo no se traduce directamente en un aumento inmediato del aprendizaje medible, su impacto en la motivación y la persistencia resulta fundamental para el desarrollo de competencias a largo plazo.

En los últimos años, el estudio del vínculo entre gamificación y flujo ha adquirido una relevancia creciente. \textcite{Oliveira_Pastushenko_Rodrigues_Toda_Palomino_Hamari_Isotani_2021} destaca, a partir de una revisión sistemática de la literatura, que el diseño gamificado puede influir significativamente en la experiencia de flujo de los usuarios, particularmente en el ámbito educativo. Los resultados muestran que la integración de elementos como recompensas, retroalimentación continua, niveles de desafío ajustables y objetivos claros favorece la motivación intrínseca y la concentración profunda. Sin embargo, también se subraya la falta de consenso metodológico en torno a qué componentes de la gamificación inducen de manera más efectiva el flujo, lo que señala la necesidad de continuar investigando este fenómeno en diferentes contextos y poblaciones.

El análisis del estado de flujo evidencia cómo la motivación intrínseca y la implicación plena del estudiante potencian la efectividad del aprendizaje. Esta comprensión sirve como punto de partida para explorar el aprendizaje adaptativo, un enfoque que ajusta dinámicamente los contenidos, la dificultad y la retroalimentación según las necesidades, habilidades y ritmo de cada estudiante. Al integrar los principios del flujo en entornos adaptativos, es posible diseñar experiencias educativas personalizadas que mantengan altos niveles de concentración y compromiso, optimizando así la adquisición de conocimientos y el desarrollo de competencias.


\section{Gamificación}

La \textbf{gamificación} ha emergido en los últimos años como una de las estrategias más prometedoras dentro de la innovación educativa, impulsada por el desarrollo de las Tecnologías de la Información y la Comunicación (TIC) y por la necesidad de adaptar los procesos de enseñanza-aprendizaje a las nuevas generaciones digitales \cite{ortiz-colon_gamificacion_2018, gallego_panoramica_nodate, jaramillo-mediavilla_impact_2024}.  
Esta metodología busca integrar en entornos no lúdicos los principios, dinámicas y elementos propios de los juegos, con el propósito de potenciar la motivación, el compromiso y el rendimiento de los estudiantes.

El término \textbf{gamificación} (también denominado ludificación) se define como el uso de mecánicas y dinámicas de juego en contextos no lúdicos \cite{ortiz-colon_gamificacion_2018}, con el fin de promover la motivación, la concentración, el esfuerzo y la fidelización.  
En esencia, consiste en trasladar la lógica de los juegos a ámbitos como la educación, la empresa o la salud, aprovechando la predisposición natural de las personas a participar, competir y superar retos \cite{gallego_panoramica_nodate}.

Desde una perspectiva más amplia, la gamificación puede entenderse como la aplicación del pensamiento del jugador y de técnicas de diseño de juegos para atraer a los usuarios, fomentar su participación y resolver problemas de manera creativa \cite{jaramillo-mediavilla_impact_2024}.  
Su finalidad no es convertir el aprendizaje en un juego, sino incorporar sus principios psicológicos —progreso, recompensa, autonomía y reto— para fortalecer el proceso educativo.

\subsection{Elementos de la gamificación}

Los fundamentos de la gamificación se estructuran en tres niveles jerárquicos —\textbf{Dinámicas}, \textbf{Mecánicas} y \textbf{Componentes}— propuestos por Werbach y Hunter (2012) \cite{ortiz-colon_gamificacion_2018,morales_elearning_noda}.  
Estos niveles permiten comprender cómo se construyen las experiencias gamificadas y cómo cada elemento contribuye al compromiso, la motivación y el aprendizaje del usuario.

\begin{table}[H]
	\centering
	\begin{tabular}{|p{3cm}|p{7cm}|p{4cm}|}
		\hline
		\textbf{Categoría} & \textbf{Descripción} & \textbf{Ejemplos en educación} \\ \hline
		\textbf{Dinámicas} & Son los elementos más abstractos de la gamificación, relacionados con las necesidades, deseos y emociones humanas que sustentan la experiencia. Representan las motivaciones profundas que impulsan la participación del usuario. & Narrativa, progresión, autonomía, sentido de pertenencia, curiosidad, altruismo, restricción temporal. \\ \hline
		\textbf{Mecánicas} & Constituyen los procesos y reglas que traducen las dinámicas en acciones concretas. Estructuran la experiencia de juego al definir cómo interactúan los participantes con el sistema. & Retos, recompensas, cooperación, competencia, retroalimentación, coleccionismo, exploración. \\ \hline
		\textbf{Componentes} & Son las manifestaciones visibles y tangibles de las mecánicas y dinámicas. Representan los elementos concretos con los que el usuario interactúa directamente. & Puntos, niveles, insignias, rankings, avatares, misiones, barras de progreso, logros. \\ \hline
	\end{tabular}
	\caption{Categorías principales de la gamificación. Fuente: Adaptado de Werbach y Hunter (2012).}
\end{table}

El diseño de experiencias gamificadas efectivas depende del equilibrio entre estos tres niveles. Las dinámicas establecen la motivación subyacente; las mecánicas definen las reglas y la estructura del desafío; y los componentes proporcionan las herramientas visuales y simbólicas que refuerzan la sensación de logro y avance \cite{puritat_enhanced_2019}.

Entre los elementos más comunes, los niveles y puntos suelen emplearse para medir el progreso del estudiante y ofrecer recompensas inmediatas, mientras que las insignias y las tablas de clasificación fomentan la competencia y el reconocimiento social \cite{puritat_enhanced_2019}.  
Otros componentes clave incluyen la retroalimentación inmediata, las metas claras, la narrativa inmersiva y las barras de progreso, todos ellos diseñados para mantener la atención y proporcionar una sensación continua de crecimiento \cite{puritat_enhanced_2019}.

La progresión constituye una de las dinámicas más esenciales, pues mantiene el interés del participante al ofrecer un sentido tangible de avance y superación. Los niveles reflejan el dominio alcanzado, mientras que la narrativa otorga coherencia al proceso, convirtiendo cada actividad en parte de una historia más amplia en la que el estudiante asume el rol de protagonista \cite{ortiz-colon_gamificacion_2018, puritat_enhanced_2019}.  

Este enfoque no solo favorece la motivación, sino que también refuerza la conexión emocional con el aprendizaje, generando un entorno donde el progreso se percibe como un logro personal.

En conjunto, las dinámicas, mecánicas y componentes conforman el núcleo del diseño gamificado. Comprender su interacción permite a los educadores desarrollar entornos de aprendizaje más atractivos, en los que la motivación, la autonomía y el sentido de logro se integran como pilares del proceso educativo.

\subsection{Motivación en la gamificación}

La motivación constituye el núcleo de la gamificación y es el principal motor que impulsa la participación y el aprendizaje activo \cite{ortiz-colon_gamificacion_2018, jaramillo-mediavilla_impact_2024}.  
Motivar no solo implica captar la atención del estudiante, sino también mantener su interés, compromiso y esfuerzo sostenido a lo largo de la actividad, generando una experiencia de aprendizaje significativa \cite{ortiz-colon_gamificacion_2018}.

En la literatura se distinguen dos tipos principales de motivación \cite{ortiz-colon_gamificacion_2018,morales_elearning_noda, puritat_enhanced_2019}:

\begin{itemize}
	\item \textbf{Motivación extrínseca:} Se origina en recompensas externas al estudiante, tales como calificaciones, insignias, reconocimientos o avances visibles en un sistema de niveles. Este tipo de motivación es útil para iniciar la participación y establecer hábitos de aprendizaje, aunque por sí sola no garantiza un compromiso duradero.  
	\item \textbf{Motivación intrínseca:} Surge del interés personal, la curiosidad y la satisfacción de realizar una actividad por el placer de aprender o por el desafío que representa. La motivación intrínseca es más sostenible y está asociada a un aprendizaje profundo y significativo \cite{ortiz-colon_gamificacion_2018}.  
\end{itemize}

Un objetivo central de la gamificación educativa es fomentar la transición de la motivación extrínseca hacia la intrínseca, de modo que los estudiantes no solo actúen por recompensas externas, sino que encuentren valor y disfrute en el propio proceso de aprendizaje \cite{sanchez-pacheco_enfoque_2020}.  

Elementos como la autonomía en la toma de decisiones, la percepción de competencia y la relevancia del contenido son determinantes para favorecer esta transformación \cite{gallego_panoramica_nodate}.  

Asimismo, la motivación puede reforzarse mediante la combinación de dinámicas, mecánicas y componentes gamificados. Por ejemplo, los desafíos progresivos, la retroalimentación inmediata, las metas claras y la narrativa inmersiva potencian la sensación de logro, dominio y pertenencia, consolidando la motivación intrínseca.  

La integración estratégica de estos elementos permite diseñar experiencias educativas que no solo retienen la atención del estudiante, sino que también promueven la iniciativa, la creatividad y la autoeficacia, estableciendo un ciclo positivo de aprendizaje continuo.

\section{Análisis de videojuegos educativos}


\subsection{Human Resource Machine}

Human Resource Machine, un videojuego de carácter educativo que enseña principios fundamentales de la programación a través de una metáfora de oficina. En este entorno, el jugador asume el papel de un trabajador encargado de ejecutar tareas mediante la construcción de algoritmos visuales que imitan las operaciones básicas del lenguaje ensamblador. A medida que el jugador avanza, se enfrenta a desafíos que exigen una planificación lógica y la optimización de procesos, fomentando así la comprensión de estructuras secuenciales, condicionales y repetitivas. Este diseño combina de forma efectiva la mecánica del rompecabezas con el aprendizaje progresivo de la programación, convirtiéndolo en un ejemplo paradigmático de cómo la simplicidad lúdica puede reforzar conceptos técnicos complejos \cite{Zhang_Wong_Chan_2023, heithausen_look_nodate}.


\subsection{CodeCombat}

Otro caso notable es CodeCombat, un juego de rol en línea que utiliza un entorno narrativo interactivo para enseñar lenguajes de programación como Python y JavaScript. Los jugadores deben escribir código real para controlar a sus personajes y superar obstáculos, lo que convierte la experiencia de juego en una práctica directa de programación aplicada. A diferencia de otros entornos educativos más teóricos, CodeCombat integra el aprendizaje dentro de la propia dinámica del juego, donde el progreso depende directamente de las decisiones lógicas del jugador. Su estructura por niveles, combinada con la retroalimentación inmediata y el componente cooperativo en línea, lo posiciona como una herramienta eficaz tanto para la enseñanza formal como para el aprendizaje autodidacta \cite{Zhang_Wong_Chan_2023}.


\subsection{Scratch}
En una línea más accesible para niveles educativos iniciales, Scratch se ha consolidado como una de las plataformas más influyentes en la enseñanza de la programación y el pensamiento computacional. Desarrollada por el MIT Media Lab, esta herramienta permite a los usuarios crear proyectos interactivos mediante bloques visuales de código, eliminando la necesidad de escribir sintaxis compleja. Su enfoque basado en la exploración, la creatividad y el aprendizaje colaborativo promueve el desarrollo de habilidades lógicas y de resolución de problemas desde edades tempranas \cite{maloney_scratch_2010}. Además, su comunidad global fomenta el intercambio de proyectos y el aprendizaje social, lo que refuerza su valor pedagógico más allá del aula tradicional \cite{velasco-ramirez_impacto_2023,maloney_scratch_2010}.


\subsection{Minerva}

Minerva es un videojuego educativo adaptativo desarrollado por investigadores de la Universidad de Ajou (Corea del Sur), orientado a la enseñanza de los fundamentos de la programación en estudiantes de educación primaria. Su propuesta combina una narrativa de ciencia ficción inmersiva con un entorno de exploración flexible, en el que el alumno controla un robot encargado de reparar una nave espacial a la deriva. A través de esta experiencia, el juego introduce de manera progresiva conceptos clave de programación como entrada y salida de datos, operaciones matemáticas, estructuras de repetición y toma de decisiones, favoreciendo la motivación y la comprensión contextual de los contenidos \cite{lindberg2017improving}.

Uno de los aportes más relevantes de Minerva es su modelo de adaptación doble, basado en el perfilado inicial del estudiante según su estilo de juego y su estilo de aprendizaje. Por un lado, el juego ajusta dinámicamente la experiencia lúdica utilizando el modelo de Bartle, modificando elementos como la cantidad de enemigos, la presencia de herramientas de exploración, los sistemas de comunicación o la retroalimentación por puntuación. Por otro lado, el contenido educativo se organiza siguiendo el modelo de Honey y Mumford, variando el orden de presentación de los bloques teóricos, instructivos, demostrativos y prácticos para adaptarse a las preferencias cognitivas de cada alumno \cite{lindberg2017improving}.

La evaluación del videojuego con estudiantes de sexto grado evidenció resultados positivos en términos de motivación, interacción social y personalización del aprendizaje. El sistema demostró ser capaz de actualizar el perfil del jugador en tiempo real, ajustando la experiencia ante cambios en su comportamiento. No obstante, también se identificaron retos, como la tendencia de algunos alumnos a omitir los tutoriales, lo que señala la necesidad de una adaptación más profunda. En este sentido, la literatura propone la incorporación de objetivos personalizados y sistemas de recompensas, consolidando a Minerva como una referencia sólida para el diseño de plataformas educativas adaptativas orientadas al aprendizaje de la programación \cite{lindberg2017improving}.

\section{Videojuegos educativos adaptativos}

Los videojuegos educativos adaptativos constituyen una evolución significativa
respecto a los videojuegos educativos tradicionales, al incorporar mecanismos
capaces de ajustar dinámicamente la experiencia de aprendizaje en función del
desempeño, las preferencias y el perfil cognitivo de cada estudiante. A diferencia
de los entornos estáticos, donde todos los usuarios reciben el mismo contenido
en el mismo orden y con la misma dificultad, los sistemas adaptativos persiguen
mantener al estudiante dentro de la zona de desarrollo próximo, evitando tanto
el aburrimiento ante tareas demasiado sencillas como la ansiedad ante desafíos
que superan sus capacidades actuales \cite{Csikszentmihalyi_2008, kiili_design_2012}.
Este principio se sustenta directamente en la teoría del flujo: el sistema debe
calibrar continuamente el nivel de reto para preservar el estado de inmersión y
compromiso sostenido que maximiza la efectividad del aprendizaje.

Desde una perspectiva conceptual, la adaptación en estos sistemas opera sobre
dos dimensiones interdependientes. La primera es la \textbf{adaptación lúdica},
que modifica elementos propios de la experiencia de juego —como la cantidad de
obstáculos, el tiempo disponible, las ayudas visuales o la complejidad de los
escenarios— con el objetivo de sostener la motivación del jugador. La segunda es
la \textbf{adaptación pedagógica}, que ajusta el contenido educativo en sí mismo:
el orden de presentación de los conceptos, la granularidad de las explicaciones,
el tipo de retroalimentación proporcionada o la selección de los ejercicios a
resolver. Ambas dimensiones pueden actuar de forma independiente o combinada,
dependiendo del diseño del sistema \cite{shum_personalised_2023, lindberg2017improving}.
El videojuego Minerva ilustra con claridad esta dualidad: por un lado, emplea el
modelo de Bartle para ajustar la experiencia lúdica según el estilo de juego del
estudiante; por otro, organiza el contenido instruccional siguiendo el modelo de
Honey y Mumford para adecuarse a sus preferencias cognitivas \cite{lindberg2017improving}.

El componente técnico que hace posible esta adaptación es el \textbf{modelo del
	estudiante}, una representación interna que el sistema construye y actualiza
de forma continua a partir de las interacciones del jugador. Este modelo recoge
variables como el historial de aciertos y errores, los tiempos de resolución,
los patrones de uso de ayudas y la frecuencia de repetición de niveles, entre
otros indicadores. A partir de esta información, un agente de adaptación —que
puede basarse en reglas heurísticas, algoritmos estadísticos o técnicas de
aprendizaje automático— toma decisiones sobre qué contenido o configuración
de dificultad ofrecer a continuación \cite{shum_personalised_2023}. La calidad
y precisión de este modelo determina en gran medida la efectividad de la
adaptación: un modelo superficial generará ajustes genéricos, mientras que uno
de alta granularidad permitirá una personalización genuina que atienda las
necesidades específicas de cada estudiante.

A pesar de sus ventajas, la literatura identifica limitaciones recurrentes en la
implementación de videojuegos educativos adaptativos. Una de las más señaladas
es la tendencia de los estudiantes a eludir los mecanismos de adaptación, por
ejemplo, omitiendo tutoriales o seleccionando estratégicamente respuestas para
manipular la dificultad percibida por el sistema \cite{lindberg2017improving}.
Asimismo, la mayoría de las soluciones existentes concentran sus esfuerzos en
la adaptación de la experiencia del estudiante, descuidando la visibilidad que
el docente necesita sobre dicho proceso. Esta brecha es precisamente la que
justifica la necesidad de integrar los mecanismos de adaptación con tableros de
analíticas de aprendizaje, de modo que el profesor pueda comprender no solo el
resultado final del estudiante, sino el proceso adaptativo que lo originó y
actuar pedagógicamente sobre él \cite{rahimi2021learning, shum_personalised_2023}.


\section{Tableros de analíticas de aprendizaje en entornos de videojuegos}

Los tableros de analíticas de aprendizaje (LAD, por sus siglas en inglés) se han consolidado como una pieza fundamental dentro del campo del Análisis de Aprendizaje (\textit{Learning Analytics}), actuando como la capa de visualización que permite cerrar el ciclo de retroalimentación educativa. Se definen como herramientas de apoyo para los docentes en el monitoreo y análisis del progreso de los estudiantes, con el objetivo de tomar decisiones informadas sobre el desempeño y la comprensión de los contenidos. Estas interfaces integran distintos tipos de gráficos y tablas que se alimentan de los datos generados por las interacciones de los estudiantes en los entornos virtuales \cite{long2011penetrating, rahimi2021learning}.

En el contexto de los videojuegos educativos, los LAD cumplen un papel particularmente relevante. A diferencia de los entornos virtuales de aprendizaje tradicionales, donde la interacción suele limitarse a la lectura o a clics, los videojuegos generan datos de alta granularidad sobre las decisiones, errores y estrategias de los jugadores. En este sentido, los LAD actúan como un puente entre la experiencia de juego (\textit{gameplay}) y el aprendizaje, haciendo visible el progreso cognitivo que, de otro modo, permanecería oculto tras la interfaz lúdica para estudiantes, docentes o padres \cite{rahimi2021learning}.

Su función principal es dotar a los profesores del contexto necesario para comprender el desempeño estudiantil, facilitando una intervención en el momento adecuado. Estos tableros rastrean todas las interacciones y el tiempo dedicado a cada tarea, ofreciendo una cantidad de información que supera con creces la que un docente podría captar de forma manual en un aula tradicional. No solo funcionan como visualizadores de datos, sino que permiten segmentar a los estudiantes para priorizar los esfuerzos pedagógicos, identificar patrones de desmotivación por falta de desafío o disparar alertas automáticas en el momento exacto en que un estudiante se enfrenta a dificultades \cite{rahimi2021learning}. 

No obstante, la literatura advierte que la mera visualización de datos no garantiza la mejora del aprendizaje. Existe el riesgo de asumir que las métricas de interacción reflejan de forma absoluta y objetiva el conocimiento del estudiante. Por ello, los LAD deben ser concebidos como herramientas de mediación que facilitan la retroalimentación formativa, permitiendo al docente ajustar su estrategia de enseñanza sobre la marcha de manera crítica \cite{ilearn2023dashboards}.

\subsection{Principios de diseño de tableros de analíticas de aprendizaje}

La literatura especializada propone un conjunto de principios para que los LAD sean pedagógicamente útiles, enfatizando que deben ir más allá de mostrar métricas de desempeño superficial y centrarse en el aprendizaje real. Para estructurar la información de manera significativa, los LAD pueden utilizar distintos marcos de referencia \cite{susnjak_learning_2022}:

\begin{itemize}
	\item \textbf{Marco social:} basado en la comparación del desempeño del estudiante con el de sus pares.
	\item \textbf{Marco de logro:} fundamentado en el cumplimiento de un objetivo específico o criterio de éxito predefinido.
	\item \textbf{Marco de progreso:} centrado en la evolución histórica del propio estudiante a lo largo del tiempo, fomentando la autorregulación.
\end{itemize}

Respecto al marco social, la literatura recomienda emplearlo con precaución. Desde una perspectiva psicológica, la comparación directa con estudiantes sobresalientes puede impactar negativamente en la autoeficacia y desmotivar a aquellos de bajo rendimiento \cite{rahimi2021learning}.

Un principio clave es que los LAD deben ser \textbf{accionables}. Esto implica que, de forma automatizada, deben ser capaces de categorizar el estado de los estudiantes (por ejemplo, ``Óptimo'', ``Necesita refuerzo'', ``En riesgo'') para que el profesor pueda tomar decisiones sin necesidad de interpretar manualmente los gráficos subyacentes \cite{rahimi2021learning}. 

Asimismo, los tableros más avanzados adoptan un enfoque \textbf{proactivo}: no solo muestran el estado actual del rendimiento, sino que también brindan un conjunto de acciones de respuesta sugeridas para mejorar la situación cuando un estudiante cae por debajo de un umbral base \cite{rahimi2021learning}. 

Finalmente, la regla de oro en la visualización de métricas es que la información esencial debe comprenderse con un solo vistazo. El dashboard debe ser fácil de interpretar tanto para alumnos como para docentes, seleccionando estrictamente los gráficos clave que aporten contexto real y evitando múltiples visualizaciones que solo generen ruido visual \cite{ilearn2023dashboards}.

\subsection{Casos de ejemplo de tableros de analíticas de aprendizaje}

\textbf{Kahoot!}

Kahoot! es una de las herramientas de gamificación más difundidas y reconocidas en el ámbito educativo contemporáneo. Se define como una plataforma de aprendizaje colaborativo y un \textit{Student Response System} (SRS), que permite a los docentes diseñar cuestionarios, encuestas y actividades de discusión con el objetivo de evaluar el conocimiento, la participación y el nivel de compromiso de los estudiantes en tiempo real \cite{singh2024exploring}.

Desde el punto de vista de sus mecánicas de juego, Kahoot! integra principios propios del diseño lúdico en un contexto educativo, incorporando elementos como sistemas de puntuación, música dinámica y una interfaz visual basada en colores llamativos. Estas características transforman el aula tradicional en un entorno similar a un concurso televisivo, donde el docente asume el rol de moderador y los estudiantes participan activamente compitiendo entre sí. El uso de límites de tiempo y tablas de clasificación refuerza el componente competitivo, incentivando respuestas rápidas y promoviendo un pensamiento ágil durante la actividad \cite{singh2024exploring}.

Diversos estudios coinciden en que el uso de Kahoot! genera beneficios pedagógicos significativos, especialmente en términos de motivación y compromiso. La plataforma favorece la motivación intrínseca del alumnado, incrementa la satisfacción durante el proceso de aprendizaje y contribuye a una mayor asistencia a clase. Asimismo, fomenta un enfoque de aprendizaje activo y centrado en el estudiante, permitiendo reforzar contenidos teóricos y prácticos de manera dinámica y participativa \cite{singh2024exploring}.

No obstante, un análisis desde la perspectiva de las analíticas revela que las capacidades de Kahoot! se centran en métricas de desempeño extrínseco, como preguntas acertadas y velocidad de respuesta. Como advierte la literatura, si bien esto incrementa la motivación situacional, un énfasis excesivo en la competencia y las recompensas puede derivar en un compromiso superficial, desplazando el foco del aprendizaje profundo hacia la simple obtención de puntos. Además, su capacidad de intervención proactiva para el docente es limitada, restringiéndose a la retroalimentación reactiva al finalizar la actividad \cite{singh2024exploring}.

\textbf{Duolingo}

Duolingo es una plataforma educativa diseñada con el objetivo de proporcionar educación de alta calidad a escala global, mediante un enfoque pedagógico que combina lecciones breves e interactivas con principios de gamificación. Su propuesta metodológica, conocida como el \textit{Método Duolingo}, se caracteriza por presentar el aprendizaje como una experiencia lúdica, sin renunciar a una base sólida en la ciencia del aprendizaje y en estándares educativos internacionales reconocidos \cite{freeman2023duolingo}.

Uno de los pilares fundamentales del Método Duolingo es el principio de \textit{aprender haciendo} (\textit{learn by doing}), que promueve la interacción directa del estudiante con el contenido desde el inicio, sin recurrir a explicaciones teóricas extensas. Este enfoque se apoya en mecanismos de aprendizaje estadístico implícito, mediante los cuales el cerebro identifica patrones a partir de la exposición repetida a estructuras cuidadosamente seleccionadas \cite{freeman2023duolingo}.

La personalización del aprendizaje constituye otro eje central de la plataforma, sustentado en el uso de inteligencia artificial. Duolingo emplea el modelo de aprendizaje automático conocido como \textit{Birdbrain}, el cual analiza diariamente un volumen masivo de interacciones para estimar el nivel de competencia del usuario y ajustar con precisión la dificultad de cada ejercicio. Este sistema aplica el concepto de zona de desarrollo próximo, ofreciendo retos que resultan estimulantes sin generar frustración, e integra técnicas de repetición espaciada para reforzar la retención a largo plazo \cite{freeman2023duolingo}.

Sin embargo, desde el punto de vista de los LAD para la intervención docente, Duolingo presenta una limitación estructural: es una plataforma diseñada fundamentalmente para el aprendizaje autodirigido. Su tablero de analíticas está más orientado a la retención del usuario (\textit{engagement}) —mediante el seguimiento de rachas y estadísticas generales— que a proveer herramientas de diagnóstico profundo que un tercero, como un profesor, pueda utilizar para mediar en el aula. Esto demuestra que una alta sofisticación algorítmica en la adaptación del contenido no siempre se traduce en dashboards accionables para la enseñanza \cite{freeman2023duolingo}.

\textbf{Physics Playground}

Physics Playground es un videojuego educativo orientado a estudiantes de secundaria o bachillerato para el aprendizaje de la física newtoniana. El objetivo principal en todos sus niveles es guiar una bola verde hasta un globo rojo. Para lograrlo, el juego propone dos modalidades de interacción:

\begin{itemize}
	\item \textbf{Niveles de dibujo:} el estudiante debe trazar artilugios sencillos (rampas, palancas, trampolines, entre otros) directamente sobre la pantalla para crear un camino físico que guíe a la bola.
	\item \textbf{Niveles de manipulación:} los alumnos interactúan con controles deslizantes (\textit{sliders}) para alterar parámetros físicos de la bola, tales como la gravedad, la elasticidad o el peso, o para aplicar fuerzas externas que permitan alcanzar el objetivo.
\end{itemize}

Este videojuego emplea un sistema de evaluación sigilosa (\textit{Stealth Assessment}) para medir el aprendizaje sin interrumpir el estado de flujo del jugador. El sistema se fundamenta en el Diseño Centrado en las Evidencias (\textit{Evidence-Centered Design}), un marco donde cada acción realizada por el jugador sirve como indicador o afirmación de su nivel de comprensión. El juego captura todas las interacciones en trazas medibles y utiliza métodos estadísticos, especialmente redes bayesianas, para realizar inferencias en tiempo real sobre el dominio del estudiante en nueve competencias clave de la física \cite{shute_stealth_2011, mislevy_evidence_2003}.

El juego incluye un tablero de analíticas denominado \textit{My Backpack}, diseñado bajo un enfoque de maestría intrapersonal. En lugar de fomentar la comparación con otros pares, este enfoque motiva al estudiante a compararse consigo mismo y con sus propias metas. Entre sus elementos se encuentran barras que muestran los niveles resueltos y las monedas recolectadas, así como gráficos que ilustran el nivel de maestría alcanzado en cada uno de los nueve conceptos fundamentales. Este caso representa una buena práctica en el diseño de LAD: utiliza modelos estadísticos robustos ocultos al usuario para traducirlos en una interfaz limpia que fomenta la reflexión y la autorregulación del aprendizaje.

\section{Selección de tecnologías}

\subsection{Base de datos}

En el desarrollo de sistemas educativos digitales, la selección del
sistema de gestión de bases de datos (SGBD) constituye un componente
crítico, ya que impacta directamente en la consistencia de la
información, el rendimiento de las consultas y la escalabilidad del
sistema. Para fundamentar la elección tecnológica, se comparan los
principales SGBD considerando criterios directamente vinculados a los
requisitos de la plataforma: la necesidad de un motor centralizado
con integridad transaccional para el backend, y un motor embebido
ligero para el videojuego educativo.

\begin{center}
	\setlength{\tabcolsep}{4pt}
	\begin{longtable}{|p{2.3cm}|p{2cm}|p{1.3cm}|p{1.5cm}|p{2cm}|p{2.5cm}|p{2cm}|}
		\caption{Comparativa de SGBD según criterios de selección.}
		\label{tab:comparacion-bd} \\
		
		\hline
		\small\textbf{SGBD} &
		\small\textbf{Modelo} &
		\small\textbf{Open Source} &
		\small\textbf{ACID} &
		\small\textbf{Despliegue} &
		\small\textbf{Escalabilidad} &
		\small\textbf{Integridad referencial} \\
		\hline
		\endfirsthead
		
		\hline
		\small\textbf{SGBD} &
		\small\textbf{Modelo} &
		\small\textbf{Open Source} &
		\small\textbf{ACID} &
		\small\textbf{Despliegue} &
		\small\textbf{Escalabilidad} &
		\small\textbf{Integridad referencial} \\
		\hline
		\endhead
		
		\hline
		\multicolumn{7}{r}{\small Continúa en la siguiente página.} \\
		\endfoot
		
		\hline
		\endlastfoot
		
		\small\textbf{PostgreSQL}~\cite{postgresql}
		& \small SQL relacional
		& \small Sí
		& \small Completo
		& \small Servidor independiente
		& \small Vertical y horizontal
		& \small Sí (foráneas y constraints) \\
		\hline
		
		\small\textbf{MySQL}~\cite{mysql}
		& \small SQL relacional
		& \small Parcial (dual)
		& \small Parcial (InnoDB)
		& \small Servidor independiente
		& \small Principalmente vertical
		& \small Sí (limitada) \\
		\hline
		
		\small\textbf{SQLite}~\cite{sqlite}
		& \small SQL embebida
		& \small Sí (dom. público)
		& \small Completo
		& \small Embebido (sin servidor)
		& \small No apta para alta concurrencia
		& \small Sí (básica) \\
		\hline
		
		\small\textbf{MongoDB}~\cite{mongodb}
		& \small NoSQL documental
		& \small Parcial (SSPL)
		& \small Parcial (v4.0)
		& \small Servidor independiente
		& \small Horizontal (sharding)
		& \small No \\
		\hline
		
		\small\textbf{Redis}~\cite{redis}
		& \small NoSQL clave-valor
		& \small Parcial (dual)
		& \small No
		& \small Servidor independiente
		& \small Horizontal
		& \small No \\
		\hline
		
	\end{longtable}
\end{center}

Del análisis comparativo se infiere que \textbf{PostgreSQL} resulta
la alternativa más adecuada para la gestión centralizada de la
información en el backend de la plataforma. Su adherencia completa al
modelo ACID (\textit{Atomicity, Consistency, Isolation, Durability})
garantiza la integridad de las operaciones críticas del sistema, como
el registro de progreso académico y la trazabilidad del desempeño
estudiantil~\cite{postgresql}. A diferencia de MySQL, cuyo
cumplimiento ACID depende del motor de almacenamiento configurado,
PostgreSQL lo garantiza de forma nativa e incondicional. Asimismo, su
soporte robusto de integridad referencial —mediante claves foráneas y
restricciones complejas— resulta indispensable para mantener la
consistencia entre las entidades del modelo de datos global: usuarios,
cursos, niveles, métricas y registros históricos. Su naturaleza
completamente open source elimina dependencias de licenciamiento,
aspecto relevante en un contexto académico donde la sostenibilidad
tecnológica a largo plazo es prioritaria.

Por su parte, \textbf{SQLite} complementa esta arquitectura al
operar como motor embebido dentro del videojuego educativo. Su
principal ventaja en este contexto no es el rendimiento en
concurrencia —para lo cual no está diseñada— sino su capacidad de
funcionar sin servidor, sin configuración de red y sin dependencias
externas, lo que la convierte en la opción natural para almacenar
localmente el estado de la sesión del jugador: niveles completados,
puntuaciones y configuraciones personalizadas~\cite{sqlite}. Al ser
de dominio público, no introduce ninguna restricción de licenciamiento
en el producto final. Estos datos son posteriormente sincronizados
con PostgreSQL en el backend, garantizando la coherencia entre la
experiencia local de juego y los registros académicos centralizados.

Las alternativas NoSQL examinadas —MongoDB y Redis— no se adecúan a
los requisitos fundamentales del sistema. La ausencia de integridad
referencial nativa en ambas y sus modelos de consistencia eventual
son incompatibles con la necesidad de trazabilidad académica precisa
que exige la plataforma. Adicionalmente, sus licencias de código
abierto parciales introducen restricciones que podrían condicionar
el uso y distribución del sistema en entornos institucionales.


\subsection{Módulo de estudiante}

El módulo del estudiante se desarrolla mediante un entorno interactivo
soportado por un motor de videojuegos. La elección del motor constituye
un aspecto clave, ya que determina las capacidades gráficas, el
rendimiento, la facilidad de integración con otros componentes del
sistema y la portabilidad del videojuego educativo. Para fundamentar
la elección, se comparan los motores más relevantes considerando
criterios directamente vinculados a los requisitos del proyecto: la
necesidad de un motor open source, exportación multiplataforma,
facilidad de integración con servicios externos y adecuación a un
contexto académico.

\begin{center}
	\setlength{\tabcolsep}{4pt}
	\begin{longtable}{|p{2cm}|p{2cm}|p{1.3cm}|p{1.8cm}|p{2cm}|p{2cm}|p{2cm}|}
		\caption{Comparativa de motores de videojuegos según criterios de selección.}
		\label{tab:motores-videojuegos} \\
		
		\hline
		\small\textbf{Motor} &
		\small\textbf{Lenguajes} &
		\small\textbf{Open Source} &
		\small\textbf{Exportación multiplataforma} &
		\small\textbf{Integración REST/ WebSocket} &
		\small\textbf{Curva de aprendizaje} &
		\small\textbf{Adecuación educativa} \\
		\hline
		\endfirsthead
		
		\hline
		\small\textbf{Motor} &
		\small\textbf{Lenguajes} &
		\small\textbf{Open Source} &
		\small\textbf{Exportación multiplataforma} &
		\small\textbf{Integración REST/ WebSocket} &
		\small\textbf{Curva de aprendizaje} &
		\small\textbf{Adecuación educativa} \\
		\hline
		\endhead
		
		\hline
		\multicolumn{7}{r}{\small Continúa en la siguiente página.} \\
		\endfoot
		
		\hline
		\endlastfoot
		
		\small\textbf{Godot 4}~\cite{godot}
		& \small GDScript, C\#, C++
		& \small Sí (MIT)
		& \small Win, Linux, Web, Android
		& \small Sí (nativa)
		& \small Baja
		& \small Alta \\
		\hline
		
		\small\textbf{Unity}~\cite{unity}
		& \small C\#
		& \small No
		& \small Win, Linux, Web, Android, iOS, consolas
		& \small Sí (via plugins)
		& \small Media
		& \small Media \\
		\hline
		
		\small\textbf{Unreal Engine}~\cite{unreal}
		& \small C++, Blueprints
		& \small No 
		& \small Win, Linux, Android, consolas
		& \small Sí (via plugins)
		& \small Alta
		& \small Baja \\
		\hline
		
		\small\textbf{Construct 3}~\cite{construct}
		& \small JavaScript
		& \small No
		& \small Web, Win (limitado)
		& \small Parcial
		& \small Muy baja
		& \small Media \\
		\hline
		
	\end{longtable}
\end{center}

Tras el análisis comparativo, se seleccionó \textbf{Godot 4} como motor principal para el desarrollo del módulo del estudiante. Es el
único motor del conjunto evaluado que combina licencia completamente
open source (MIT), exportación nativa a las plataformas objetivo del
proyecto —Windows, Android y Web— e integración nativa con servicios
externos mediante REST y WebSocket, sin depender de plugins de
terceros~\cite{godot}. Unity, aunque más maduro en términos de
ecosistema, introduce una dependencia de licenciamiento propietario
que condiciona su uso en contextos académicos y restringe la
auditoría del entorno de desarrollo. Unreal Engine, por su parte,
queda descartado por su elevada curva de aprendizaje y sus requisitos
de hardware, incompatibles con la diversidad de dispositivos presente
en entornos educativos. Construct 3, si bien ofrece la curva de
aprendizaje más baja, presenta una exportación multiplataforma
limitada y una escalabilidad insuficiente para los requisitos de
integración y adaptación dinámica de la plataforma.

Desde el punto de vista de ingeniería, Godot ofrece una arquitectura
basada en \textbf{escenas y nodos} que favorece el diseño modular:
los niveles, personajes y mecánicas pueden implementarse como
componentes reutilizables, facilitando el mantenimiento, las pruebas
y la evolución incremental del videojuego. Su sistema de señales y
las herramientas de edición en tiempo real priorizan la iteración
rápida, lo cual resulta especialmente valioso para ajustar la
dificultad, la retroalimentación y el ritmo de progresión del
aprendizaje sin interrumpir el ciclo de desarrollo.

En conjunto, Godot 4 ofrece el equilibrio más adecuado entre
\textbf{accesibilidad, portabilidad e integración} para un módulo
del estudiante centrado en experiencias interactivas y en el
seguimiento sistemático del aprendizaje.

\subsection{Módulo de profesor}

El módulo del profesor está concebido como una interfaz web que permite
visualizar, gestionar y analizar el progreso de los estudiantes dentro
de la plataforma educativa. Para su desarrollo, se requiere un
\textit{framework} frontend que ofrezca un equilibrio entre rendimiento,
escalabilidad y facilidad de integración con el back-end. Para
fundamentar la elección, se comparan los principales \textit{frameworks}
considerando criterios directamente vinculados a los requisitos del
módulo: soporte de renderizado en servidor, compatibilidad con
TypeScript, integración con APIs REST y madurez del ecosistema para
el desarrollo de dashboards analíticos.

\begin{center}
	\setlength{\tabcolsep}{4pt}
	\begin{longtable}{|p{2cm}|p{1.8cm}|p{1.3cm}|p{1.5cm}|p{1.8cm}|p{2cm}|p{2.2cm}|}
		\caption{Comparativa de \textit{frameworks} frontend según criterios de selección.}
		\label{tab:frameworks-frontend} \\
		
		\hline
		\small\textbf{Framework} &
		\small\textbf{Lenguaje} &
		\small\textbf{Open Source} &
		\small\textbf{SSR nativo} &
		\small\textbf{TypeScript} &
		\small\textbf{Integración REST} &
		\small\textbf{Madurez del ecosistema} \\
		\hline
		\endfirsthead
		
		\hline
		\small\textbf{Framework} &
		\small\textbf{Lenguaje} &
		\small\textbf{Open Source} &
		\small\textbf{SSR nativo} &
		\small\textbf{TypeScript} &
		\small\textbf{Integración REST} &
		\small\textbf{Madurez del ecosistema} \\
		\hline
		\endhead
		
		\hline
		\multicolumn{7}{r}{\small Continúa en la siguiente página.} \\
		\endfoot
		
		\hline
		\endlastfoot
		
		\small\textbf{Next.js}~\cite{nextjs}
		& \small JS / TS (React)
		& \small Sí (MIT)
		& \small Sí (SSR y SSG)
		& \small Nativo
		& \small Nativa
		& \small Alta \\
		\hline
		
		\small\textbf{React}~\cite{react}
		& \small JS / TS
		& \small Sí (MIT)
		& \small No (requiere config.)
		& \small Nativo
		& \small Via librerías
		& \small Muy alta \\
		\hline
		
		\small\textbf{Angular}~\cite{angular}
		& \small TypeScript
		& \small Sí (MIT)
		& \small Parcial (Angular Universal)
		& \small Nativo
		& \small Nativa (HttpClient)
		& \small Alta \\
		\hline
		
		\small\textbf{Vue.js}~\cite{vuejs}
		& \small JS / TS
		& \small Sí (MIT)
		& \small Limitado (Nuxt.js)
		& \small Parcial
		& \small Via librerías
		& \small Media \\
		\hline
		
		\small\textbf{SvelteKit}~\cite{sveltekit}
		& \small JS / TS
		& \small Sí (MIT)
		& \small Sí
		& \small Parcial
		& \small Via librerías
		& \small Baja \\
		\hline
		
	\end{longtable}
\end{center}

Tras el análisis comparativo, se seleccionó \textbf{Next.js} como
\textit{framework} principal para el desarrollo del módulo del profesor.
Es la opción que mejor satisface de forma simultánea los criterios
determinantes para este módulo: SSR nativo sin configuración adicional,
soporte completo de TypeScript, integración directa con APIs REST y un
ecosistema maduro con librerías especializadas en visualización de datos
como \textit{Recharts} o \textit{Chart.js}, indispensables para los
dashboards analíticos del docente~\cite{nextjs}. React, aunque comparte
la misma base tecnológica, no incluye SSR por defecto y requiere
configuración adicional para alcanzar las capacidades que Next.js
ofrece de forma nativa, lo que lo convierte en una capa inferior sobre
la que Next.js precisamente se construye. Angular, si bien presenta
integración REST nativa y soporte sólido de TypeScript, introduce una
curva de aprendizaje pronunciada y una complejidad estructural que
excede los requisitos de un módulo de tamaño medio como el presente.
Vue.js y SvelteKit quedan descartados por su soporte limitado de SSR
nativo y la menor madurez de sus ecosistemas, lo que representa un
riesgo para la mantenibilidad a largo plazo del sistema.

Su soporte nativo para rutas dinámicas y \textit{fetching} de datos
asíncronos permite la visualización eficiente de información académica
en tiempo real —progreso del estudiante, estadísticas de rendimiento
y resultados de evaluaciones— sin comprometer la experiencia de usuario.
El tipado estático que provee TypeScript refuerza adicionalmente la
mantenibilidad del código y facilita la detección temprana de errores
durante el desarrollo.




\subsection{Back-end}

El \textit{back-end} constituye el núcleo lógico de la plataforma:
gestiona la comunicación entre los componentes del sistema, procesa
la información proveniente del videojuego y del módulo del profesor,
y garantiza la integridad de los datos almacenados. Para fundamentar
la elección, se comparan los principales \textit{frameworks} de
desarrollo backend considerando criterios directamente vinculados a
los requisitos del sistema: soporte asíncrono nativo, validación
automática de datos, generación de documentación, compatibilidad con
Python —lenguaje del ecosistema tecnológico seleccionado— y madurez
de la comunidad.

\begin{center}
	\setlength{\tabcolsep}{4pt}
	\begin{longtable}{|p{2cm}|p{1.5cm}|p{1.3cm}|p{1.8cm}|p{1.8cm}|p{1.8cm}|p{2cm}|}
		\caption{Comparativa de \textit{frameworks} backend según criterios de selección.}
		\label{tab:frameworks-backend} \\
		
		\hline
		\small\textbf{Framework} &
		\small\textbf{Lenguaje} &
		\small\textbf{Open Source} &
		\small\textbf{Soporte asíncrono} &
		\small\textbf{Validación de datos} &
		\small\textbf{Documentación automática} &
		\small\textbf{Madurez de la comunidad} \\
		\hline
		\endfirsthead
		
		\hline
		\small\textbf{Framework} &
		\small\textbf{Lenguaje} &
		\small\textbf{Open Source} &
		\small\textbf{Soporte asíncrono} &
		\small\textbf{Validación de datos} &
		\small\textbf{Documentación automática} &
		\small\textbf{Madurez de la comunidad} \\
		\hline
		\endhead
		
		\hline
		\multicolumn{7}{r}{\small Continúa en la siguiente página.} \\
		\endfoot
		
		\hline
		\endlastfoot
		
		\small\textbf{FastAPI}~\cite{fastapi}
		& \small Python
		& \small Sí (MIT)
		& \small Nativo (ASGI)
		& \small Nativa (Pydantic)
		& \small Nativa (OpenAPI)
		& \small Media (en crecimiento) \\
		\hline
		
		\small\textbf{Django}~\cite{django}
		& \small Python
		& \small Sí (BSD)
		& \small Limitado (WSGI)
		& \small Via formularios
		& \small No
		& \small Muy alta \\
		\hline
		
		\small\textbf{Flask}~\cite{flask}
		& \small Python
		& \small Sí (BSD)
		& \small Parcial (extensiones)
		& \small Manual
		& \small No
		& \small Alta \\
		\hline
		
		\small\textbf{Express.js}~\cite{express}
		& \small JavaScript
		& \small Sí (MIT)
		& \small Nativo (Node.js)
		& \small Manual
		& \small No
		& \small Muy alta \\
		\hline
		
		\small\textbf{Laravel}~\cite{laravel}
		& \small PHP
		& \small Sí (MIT)
		& \small Limitado
		& \small Via validadores
		& \small No
		& \small Alta \\
		\hline
		
		\small\textbf{Spring Boot}~\cite{springboot}
		& \small Java
		& \small Sí (Apache)
		& \small Sí (WebFlux)
		& \small Via anotaciones
		& \small Parcial (Swagger)
		& \small Muy alta \\
		\hline
		
	\end{longtable}
\end{center}

Tras el análisis comparativo, se seleccionó \textbf{FastAPI} como
\textit{framework} back-end para el desarrollo del sistema. Es la
única opción del conjunto evaluado que combina de forma nativa los
tres criterios más determinantes para la plataforma: soporte
asíncrono completo mediante ASGI, validación automática de datos
con Pydantic y generación de documentación interactiva mediante
OpenAPI~\cite{fastapi}. Django y Flask, aunque comparten el
ecosistema Python, presentan limitaciones estructurales relevantes:
Django fue diseñado sobre el estándar WSGI, lo que restringe su
capacidad para manejar concurrencia real, y Flask delega la
validación de datos a configuraciones manuales, incrementando la
superficie de error en un sistema que intercambia métricas académicas
entre múltiples clientes. Express.js, si bien ofrece soporte
asíncrono nativo, rompe la homogeneidad del ecosistema Python del
proyecto e introduce una gestión manual de tipos y validaciones
incompatible con los requisitos de consistencia de datos del sistema.
Laravel y Spring Boot quedan descartados por introducir lenguajes
ajenos al ecosistema tecnológico seleccionado y por una complejidad
de configuración desproporcionada respecto al alcance del proyecto.

Su integración nativa con PostgreSQL y su compatibilidad directa con
el videojuego en Godot 4 —mediante peticiones REST— y con el módulo
del profesor en Next.js garantizan una comunicación rápida, segura y
eficiente entre los tres componentes de la plataforma. El soporte
para programación asíncrona permite además manejar eventos en tiempo
real, como el registro de resultados y la adaptación dinámica de la
dificultad, sin comprometer la estabilidad del sistema.

\section{Conclusiones parciales}